% ============================================================================
%  Cabecera del manuscrito: título, autores, afiliaciones, fecha
%  El abstract y las keywords viven en sections/00-abstract.tex
%
%  Reglas A&A (manuscript header):
%    - Separar cada autor con \and.
%    - \corrauth{correo} marca al autor de correspondencia; aa.cls v9.4 lo
%      convierte solo en una nota al pie. Un único \corrauth por artículo.
%    - \email{correo} NO se imprime: lo extrae el sistema editorial de EDPS
%      para que cada coautor autentique su ORCID.
%    - NO poner ORCIDs junto a los nombres. Se eliminan en producción.
%    - El título puede partirse con \\ y llevar \thanks (separados con \fnmsep).
% ============================================================================

\title{Título del artículo\thanks{Nota al pie del título, si hace falta:
       agencia de financiación, número de proyecto, etc.}}

% Quitar por completo si no hay subtítulo.
\subtitle{Subtítulo}

\author{
  Luis Daniel Díaz\inst{1}\corrauth{correo@udea.edu.co}
  \and Nombre Coautor\inst{1}\email{coautor@udea.edu.co}
  \and Nombre Coautor\inst{2}\email{coautor@otra.edu}
}

\institute{
  Instituto de Física, Universidad de Antioquia, Calle 70 No. 52-21,
  Medellín, Colombia
  \and
  Segunda afiliación, Ciudad, País
}

\date{Received DD Month YYYY / Accepted DD Month YYYY}
